\section{Examples: Computing Character Tables by Orthogonality Relations}

In this section we illustrate how to compute character tables using the row
and column orthogonality relations.

Recall that if
\[
    \operatorname{Irr}(G)=\{\chi_1,\dots,\chi_r\}
\]
is the set of irreducible complex characters of \(G\), and if
\[
    C_1,\dots,C_r
\]
are the conjugacy classes of \(G\), with representatives \(g_1,\dots,g_r\),
then the row orthogonality relations say that
\[
    \langle \chi_i,\chi_j\rangle
    =
    \frac{1}{|G|}
    \sum_{\ell=1}^r
    |C_\ell|\chi_i(g_\ell)\overline{\chi_j(g_\ell)}
    =
    \delta_{ij}.
\]
The column orthogonality relations say that
\[
    \sum_{\chi\in \operatorname{Irr}(G)}
    \chi(g_i)\overline{\chi(g_j)}
    =
    \begin{cases}
        |C_G(g_i)|, & g_i\text{ and }g_j\text{ are conjugate},\\
        0, & g_i\text{ and }g_j\text{ are not conjugate}.
    \end{cases}
\]
In particular, for a representative \(g_i\),
\[
    \sum_{\chi\in \operatorname{Irr}(G)}
    |\chi(g_i)|^2
    =
    |C_G(g_i)|.
\]

\begin{example}[The character table of \(S_3\)]
    The conjugacy classes of \(S_3\) are represented by
    \[
        1,\qquad (12),\qquad (123).
    \]
    Their sizes and centralizer orders are
    \[
        \begin{array}{c|ccc}
            \text{representative} & 1 & (12) & (123)\\
            \hline
            \text{class size} & 1 & 3 & 2\\
            |C_{S_3}(g)| & 6 & 2 & 3
        \end{array}
    \]
    Hence \(S_3\) has exactly three irreducible complex characters.

    There are two one-dimensional characters:
    \[
        \chi_{\mathrm{triv}}=(1,1,1),
        \qquad
        \chi_{\mathrm{sgn}}=(1,-1,1).
    \]
    If the remaining irreducible character has degree \(d\), then
    \[
        1^2+1^2+d^2=|S_3|=6,
    \]
    so
    \[
        d=2.
    \]
    Write the remaining character as
    \[
        \chi=(2,a,b),
    \]
    where \(a=\chi((12))\) and \(b=\chi((123))\).

    We first use the column orthogonality relation. For the transposition class,
    \[
        |\chi_{\mathrm{triv}}((12))|^2
        +
        |\chi_{\mathrm{sgn}}((12))|^2
        +
        |\chi((12))|^2
        =
        |C_{S_3}((12))|.
    \]
    Thus
    \[
        1^2+(-1)^2+|a|^2=2,
    \]
    and hence
    \[
        a=0.
    \]
    For the \(3\)-cycle class,
    \[
        1^2+1^2+|b|^2=3,
    \]
    so
    \[
        |b|=1.
    \]

    To determine the sign of \(b\), use row orthogonality with the trivial character:
    \[
        0=\langle \chi,\chi_{\mathrm{triv}}\rangle
        =
        \frac{1}{6}
        \bigl(
            1\cdot 2
            +
            3\cdot a
            +
            2\cdot b
        \bigr).
    \]
    Since \(a=0\), this becomes
    \[
        2+2b=0,
    \]
    hence
    \[
        b=-1.
    \]

    Therefore the character table of \(S_3\) is
    \[
        \begin{array}{c|ccc}
            S_3 & 1 & (12) & (123)\\
            \hline
            \text{class size} & 1 & 3 & 2\\
            \hline
            \chi_{\mathrm{triv}} & 1 & 1 & 1\\
            \chi_{\mathrm{sgn}} & 1 & -1 & 1\\
            \chi & 2 & 0 & -1
        \end{array}
    \]
\end{example}

\begin{example}[The character table of \(D_8\)]
    Let
    \[
        D_8=\langle r,s\mid r^4=s^2=1,\ srs=r^{-1}\rangle
    \]
    be the dihedral group of order \(8\). Its conjugacy classes are represented by
    \[
        1,\qquad r^2,\qquad r,\qquad s,\qquad rs.
    \]
    More explicitly,
    \[
        \{1\},\qquad
        \{r^2\},\qquad
        \{r,r^3\},\qquad
        \{s,r^2s\},\qquad
        \{rs,r^3s\}.
    \]
    Their sizes and centralizer orders are
    \[
        \begin{array}{c|ccccc}
            \text{representative} & 1 & r^2 & r & s & rs\\
            \hline
            \text{class size} & 1 & 1 & 2 & 2 & 2\\
            |C_{D_8}(g)| & 8 & 8 & 4 & 4 & 4
        \end{array}
    \]
    Hence \(D_8\) has exactly five irreducible complex characters.

    Since
    \[
        D_8/[D_8,D_8]\cong C_2\times C_2,
    \]
    there are four one-dimensional characters. They are determined by the
    possible values of \(r\) and \(s\), both of which must be \(\pm 1\):
    \[
        \begin{array}{c|ccccc}
            & 1 & r^2 & r & s & rs\\
            \hline
            \chi_{++} & 1 & 1 & 1 & 1 & 1\\
            \chi_{+-} & 1 & 1 & 1 & -1 & -1\\
            \chi_{-+} & 1 & 1 & -1 & 1 & -1\\
            \chi_{--} & 1 & 1 & -1 & -1 & 1
        \end{array}
    \]
    Let the remaining irreducible character be
    \[
        \chi=(2,a,b,c,d),
    \]
    where the entries correspond to the classes
    \[
        1,\quad r^2,\quad r,\quad s,\quad rs.
    \]
    Its degree is \(2\), because
    \[
        1^2+1^2+1^2+1^2+\chi(1)^2=|D_8|=8.
    \]

    We now use column orthogonality. For the class of \(r\),
    \[
        1^2+1^2+(-1)^2+(-1)^2+|b|^2
        =
        |C_{D_8}(r)|
        =
        4.
    \]
    Thus
    \[
        b=0.
    \]
    Similarly, for the two reflection classes,
    \[
        c=0,
        \qquad
        d=0.
    \]
    For the central element \(r^2\), column orthogonality gives
    \[
        1^2+1^2+1^2+1^2+|a|^2
        =
        |C_{D_8}(r^2)|
        =
        8.
    \]
    Hence
    \[
        |a|=2.
    \]

    To determine \(a\), use row orthogonality with the trivial character:
    \[
        0
        =
        \langle \chi,\chi_{++}\rangle
        =
        \frac{1}{8}
        \bigl(
            1\cdot 2
            +
            1\cdot a
            +
            2\cdot b
            +
            2\cdot c
            +
            2\cdot d
        \bigr).
    \]
    Since \(b=c=d=0\), this gives
    \[
        2+a=0,
    \]
    so
    \[
        a=-2.
    \]

    Therefore the character table of \(D_8\) is
    \[
        \begin{array}{c|ccccc}
            D_8 & 1 & r^2 & r & s & rs\\
            \hline
            \text{class size} & 1 & 1 & 2 & 2 & 2\\
            \hline
            \chi_{++} & 1 & 1 & 1 & 1 & 1\\
            \chi_{+-} & 1 & 1 & 1 & -1 & -1\\
            \chi_{-+} & 1 & 1 & -1 & 1 & -1\\
            \chi_{--} & 1 & 1 & -1 & -1 & 1\\
            \chi & 2 & -2 & 0 & 0 & 0
        \end{array}
    \]
\end{example}

\begin{example}[The character table of \(S_4\)]
    The conjugacy classes of \(S_4\) are determined by cycle type. We take
    representatives
    \[
        1,\qquad
        (12),\qquad
        (12)(34),\qquad
        (123),\qquad
        (1234).
    \]
    The class sizes and centralizer orders are
    \[
        \begin{array}{c|ccccc}
            \text{representative}
            & 1 & (12) & (12)(34) & (123) & (1234)\\
            \hline
            \text{class size}
            & 1 & 6 & 3 & 8 & 6\\
            |C_{S_4}(g)|
            & 24 & 4 & 8 & 3 & 4
        \end{array}
    \]
    Hence \(S_4\) has exactly five irreducible complex characters.

    We already know two one-dimensional characters:
    \[
        \chi_{\mathrm{triv}}
        =
        (1,1,1,1,1),
    \]
    and
    \[
        \chi_{\mathrm{sgn}}
        =
        (1,-1,1,1,-1).
    \]

    Next consider the permutation representation of \(S_4\) on
    \[
        \mathbb{C}^4.
    \]
    Its character is the number of fixed points of a permutation. Hence its
    values are
    \[
        (4,2,0,1,0).
    \]
    Removing the one-dimensional trivial subrepresentation, we get the standard
    representation \(V_{\mathrm{std}}\), whose character is
    \[
        \chi_{\mathrm{std}}
        =
        (4,2,0,1,0)-(1,1,1,1,1)
        =
        (3,1,-1,0,-1).
    \]
    We check that this character is irreducible:
    \[
        \langle \chi_{\mathrm{std}},\chi_{\mathrm{std}}\rangle
        =
        \frac{1}{24}
        \left(
            1\cdot 3^2
            +
            6\cdot 1^2
            +
            3\cdot (-1)^2
            +
            8\cdot 0^2
            +
            6\cdot (-1)^2
        \right)
        =
        1.
    \]
    Therefore \(\chi_{\mathrm{std}}\) is irreducible.

    Twisting by the sign character gives another irreducible character:
    \[
        \chi_{\mathrm{sgn}}\chi_{\mathrm{std}}
        =
        (3,-1,-1,0,1).
    \]

    So far we have four irreducible characters, of degrees
    \[
        1,\quad 1,\quad 3,\quad 3.
    \]
    Let the remaining irreducible character be
    \[
        \chi=(2,a,b,c,d),
    \]
    where the entries correspond to
    \[
        1,\quad (12),\quad (12)(34),\quad (123),\quad (1234).
    \]
    Its degree is \(2\), since
    \[
        1^2+1^2+3^2+3^2+\chi(1)^2
        =
        |S_4|
        =
        24.
    \]

    We first use column orthogonality. For the transposition class \((12)\),
    \[
        1^2+(-1)^2+1^2+(-1)^2+|a|^2
        =
        |C_{S_4}((12))|
        =
        4.
    \]
    Therefore
    \[
        a=0.
    \]
    For the \(4\)-cycle class \((1234)\),
    \[
        1^2+(-1)^2+(-1)^2+1^2+|d|^2
        =
        |C_{S_4}((1234))|
        =
        4,
    \]
    so
    \[
        d=0.
    \]
    For the double transposition class \((12)(34)\),
    \[
        1^2+1^2+(-1)^2+(-1)^2+|b|^2
        =
        |C_{S_4}((12)(34))|
        =
        8,
    \]
    hence
    \[
        |b|=2.
    \]
    For the \(3\)-cycle class \((123)\),
    \[
        1^2+1^2+0^2+0^2+|c|^2
        =
        |C_{S_4}((123))|
        =
        3,
    \]
    hence
    \[
        |c|=1.
    \]

    Column orthogonality has determined
    \[
        a=0,\qquad d=0,\qquad |b|=2,\qquad |c|=1.
    \]
    To determine the signs of \(b\) and \(c\), we now use row orthogonality.

    Orthogonality with \(\chi_{\mathrm{std}}\) gives
    \[
        0
        =
        \langle \chi,\chi_{\mathrm{std}}\rangle
        =
        \frac{1}{24}
        \left(
            1\cdot 2\cdot 3
            +
            6\cdot a\cdot 1
            +
            3\cdot b\cdot (-1)
            +
            8\cdot c\cdot 0
            +
            6\cdot d\cdot (-1)
        \right).
    \]
    Since \(a=d=0\), this becomes
    \[
        6-3b=0.
    \]
    Hence
    \[
        b=2.
    \]
    Orthogonality with the trivial character gives
    \[
        0
        =
        \langle \chi,\chi_{\mathrm{triv}}\rangle
        =
        \frac{1}{24}
        \left(
            1\cdot 2
            +
            6\cdot a
            +
            3\cdot b
            +
            8\cdot c
            +
            6\cdot d
        \right).
    \]
    Using \(a=d=0\) and \(b=2\), we obtain
    \[
        2+6+8c=0.
    \]
    Therefore
    \[
        c=-1.
    \]

    Thus the character table of \(S_4\) is
    \[
        \begin{array}{c|ccccc}
            S_4
            & 1 & (12) & (12)(34) & (123) & (1234)\\
            \hline
            \text{class size}
            & 1 & 6 & 3 & 8 & 6\\
            \hline
            \chi_{\mathrm{triv}}
            & 1 & 1 & 1 & 1 & 1\\
            \chi_{\mathrm{sgn}}
            & 1 & -1 & 1 & 1 & -1\\
            \chi_{\mathrm{std}}
            & 3 & 1 & -1 & 0 & -1\\
            \chi_{\mathrm{sgn}}\chi_{\mathrm{std}}
            & 3 & -1 & -1 & 0 & 1\\
            \chi
            & 2 & 0 & 2 & -1 & 0
        \end{array}
    \]
\end{example}
